LLM-generated code that passes dense differential test suites --- 764 tests per task --- routinely violates formal contracts on the same inputs: programs returning the ground-truth output on contract-violating inputs nonetheless fail reference postconditions 40\% of the time. We introduce an oracle isolation design that separates contract oracle semantic strength from input generation strategy, using ContractEval's contract-violating test inputs (CVTs) as a controlled fixed input set to directly compare differential and contract oracles on identical inputs. Across 364 HumanEval+/MBPP+ tasks and 3 LLM families (Experiment A), we find an oracle-isolation gap of 0.40 (4$\times$ the predicted threshold; Wilcoxon $p = 5.88 \times 10^{-38}$) and a 40\% contract-unique mass on CVT inputs --- programs output-equal to ground truth but violating reference contracts on these contract-relevant inputs. Adaptive property-based testing via icontract-hypothesis adds an additional $\sim$10 percentage points of violations beyond static oracle inputs ($p = 2.64 \times 10^{-22}$), consistently across all 5 tested model families (Experiment B), confirming that oracle semantics and adaptive search contribute independent detection layers. Our findings show that contract oracles and differential test oracles are qualitatively different evaluation instruments, arguing for oracle-type diversity as a first-class dimension in LLM code evaluation.
